\section{Finite actions, sections, and rotation data}
\label{sec:partB}

In this section $U=X\setminus S$ is an orientable connected surface of finite type,
unless another hypothesis is explicitly stated. A universal covering
$p:\widetilde U\to U$ is fixed, and $\pi=\operatorname{Deck}(p)$ acts on the left.
For a faithful action of a finite group $G$ on $U$, we write $E$ for the group of
all lifts of its elements. Thus $1\to\pi\to E\to G\to1$ is an extension of
\emph{discrete} groups. A partial section means a finite subgroup of $E$ projecting
isomorphically onto a subgroup of $G$; its equivalence relation throughout this
section is conjugacy by $\pi$. We use $G^+$ for the orientation-preserving subgroup.
The claims below also specify the corrections required when $G^+=1$.

\subsection{The plane lemma and the section correspondence}

\begin{claim}{B.1}{Proved}{Finite plane actions and simultaneous local models}
Every faithful action of a finite group $H$ on $\R^2$ by homeomorphisms is
conjugate to an orthogonal action. When $H^+\ne1$, $H^+$ is cyclic, $H$ is cyclic
or dihedral, and $(\R^2)^H$ is a single point. A nontrivial orientation-reversing
involution has a fixed line. Every stabilizer in a finite action on a surface
with universal cover $\R^2$ has an invariant disc on which the whole stabilizer
acts by an orthogonal action. These conclusions supply the two local assertions
used in \S15, folios 70--74, scans 141--149, and \S33, scan 440.
\end{claim}
\begin{proof}
We use precisely Kolev's Theorem 1.2: a compact subgroup of
$\operatorname{Homeo}(S^2)$ is conjugate to a closed subgroup of $O(3)$
\cite{bKolev2006}. Each homeomorphism of $\R^2$ extends to its one-point
compactification, and the finite group fixes the added point $\infty$.
After the stated conjugacy the group fixes a point $q$ of the round sphere.
Its action on $q^\perp$ gives an injective homomorphism into $O(2)$.
Stereographic projection from $q$ is equivariant and identifies the action on
$S^2\setminus\{q\}$ with that linear action. Finite subgroups of $SO(2)$ are the
cyclic groups of roots of unity: an element of least positive angle generates
the group, by division with remainder of angles. An orientation-reversing
orthogonal matrix in dimension two is a reflection and conjugates every rotation
to its inverse. This proves the cyclic/dihedral assertion and all assertions about
fixed sets.

For the local statement, let $H$ stabilize $x\in U$ and choose
$\widetilde x\in p^{-1}(x)$. For each $h\in H$ there is a unique lift fixing
$\widetilde x$. Uniqueness makes the assignment a homomorphism, and it is
injective because its projection to $U$ is the original action. The preceding
plane argument linearizes these lifts simultaneously. Pull back a sufficiently
small round disc about the common fixed point; choose it inside an evenly
covered neighbourhood of $\widetilde x$. Its projection is the required
$H$-invariant disc. The same assertion for a sphere follows directly from
Kolev's theorem. In particular the fixed locus of an anti-involution is locally
a diameter, with the involution interchanging its two sides. No differentiability
of the original action has entered the argument. The single-element versions
are Constantin--Kolev, Lemma 2.5, Theorems 3.1 and 4.1, and Corollary 4.2
\cite{bConstantinKolev1994}.
\end{proof}

\begin{claim}{B.2}{Refuted}{The omitted orientation hypothesis in scans 146--147}
The displayed lemma on scan 147, taken without its preceding orientation
hypothesis, is false. Its correct form is: an orientation-preserving finite-order
homeomorphism $u$ of $\C$ commuting with $T(z)=z+1$ and satisfying $u(0)=0$ is
the identity. The quotient $\C/\Z$ provides a complete proof along the route
proposed in the manuscript.
\end{claim}
\begin{proof}
Complex conjugation is a counterexample to the unqualified sentence: it has
order two, fixes zero, and commutes with $T$. For the corrected statement,
$u(k)=k$ for every integer $k$. B.1 gives at most one fixed point for a
nonidentity orientation-preserving periodic plane homeomorphism, so $u=1$.

Here is the quotient argument. Commutation with $T$ gives a finite-order
homeomorphism $\bar u$ of the annulus $A=\C/\Z$, inducing $+1$ on
$\pi_1(A)=\Z$. It preserves the orientation of $A$. The two end loops, oriented
as boundary loops of a compact subannulus, represent opposite homology classes.
Consequently an orientation-preserving map inducing $+1$ on $H_1(A;\Z)$ fixes
each end. Filling the ends extends $\bar u$ to an orientation-preserving periodic
homeomorphism of $S^2$ fixing both added points and $[0]$. By
Constantin--Kolev, Theorem 4.1, a nonidentity such map has exactly two fixed
points. Hence $\bar u=1$. Its lift $u$ is a deck transformation $T^k$; $u(0)=0$
forces $k=0$. On scan 146 the equation $G=G^\circ$ is exactly what makes the
lifted action orientation preserving. The mathematical repair therefore restores
an inherited hypothesis; it does not justify a spectral-sequence argument.
\end{proof}

\begin{claim}{B.3}{Proved}{Sections and isolated fixed points}
Let $\widetilde U\cong\R^2$. A faithful finite action on $U$ with $G^+\ne1$
has a bijection
\[
 U^G\longrightarrow \operatorname{Sec}(E\to G)/\pi.
\]
For every section subgroup $K$, $C_\pi(K)=1$. Existence of a section forces $G$
to be cyclic or dihedral. More generally, every nontrivial finite subgroup of
$E$ whose orientation-preserving part is nontrivial is contained in exactly one
maximal finite subgroup of $E$. For an orientation-preserving action the map
from points with nontrivial stabilizer to maximal partial-section classes is a
$G$-equivariant bijection. These are the theorem and corollary of \S15, folios
73--75, scans 148--151, with their stated orientation hypotheses.
\end{claim}
\begin{proof}
For $x\in U^G$, the lifts of $G$ fixing a chosen point over $x$ form a section.
Changing that point conjugates the section by the corresponding element of
$\pi$. Conversely a section $K$ is a finite faithful group of plane
homeomorphisms. B.1 supplies a unique fixed point $\widetilde x$. Its image is
fixed by $G$, and the section fixing $\widetilde x$ is $K$. These constructions
are inverse: conjugate sections have deck-translated fixed points and hence
the same image in $U$. An element of $C_\pi(K)$ preserves the unique fixed point.
A deck transformation with a fixed point is the identity. This proves the
centralizer assertion. B.1 gives the possible finite groups.

Let $K\le E$ have nontrivial positive part and fix $\widetilde x$. The stabilizer
$E_{\widetilde x}$ is finite, since it intersects $\pi$ trivially and injects into
the finite group $G$. Any finite subgroup containing $K$ has a fixed point by
B.1; that point must be $\widetilde x$. It is therefore contained in
$E_{\widetilde x}$. This proves uniqueness as an actual subgroup, before any
conjugacy quotient. In the orientation-preserving case every nontrivial partial
section has this property. Applying the same construction to each stabilizer
$G_x$ gives the final bijection and its equivariance. The argument also applies
to a nonorientable $U$ whenever the finite lifted group has nontrivial positive
part; in particular the cardinality condition $|G|\ge3$ on scan 149 ensures it.
\end{proof}

\begin{claim}{B.4}{Proved}{Reflection sections, ovals, and punctured ovals}
Let $\sigma$ be an orientation-reversing involution of $U$ with
$\widetilde U\cong\R^2$, and let $E\to C_2$ be its lift extension. There is a
bijection
\[
 \pi_0(U^\sigma)\longrightarrow\operatorname{Sec}(E\to C_2)/\pi.
\]
For the section $s$ associated with a component $Z$, inclusion induces an
isomorphism $\pi_1(Z)\cong C_\pi(s)$. Thus the centralizer is $1$ for a line and
$\Z$ for an oval. This includes punctures fixed or permuted by $\sigma$ and proves
\S33, folios 222--225, scans 440--445, in the required noncompact cases.
\end{claim}
\begin{proof}
Each order-two lift $s$ is a reflection of the plane by B.1. Its fixed set
$L$ is a line. The restriction $p:L\to U^\sigma$ is locally a homeomorphism
onto the real locus, by the local model in B.1. Its image is the entire component
$Z$ containing any of its points. To verify this without a properness assumption,
lift a path in $Z$ starting at $p(\widetilde x)$: applying $s$ gives another lift
with the same initial point, so uniqueness of path lifting fixes the lifted path.
The same argument for a homotopy shows that $L\to Z$ is its universal covering.
Every section has such a fixed line, so every section class arises. Two fixed
lines projecting to the same component contain lifts of the same point; a deck
transformation taking one lift to the other conjugates the associated involutions.
The two lifts fixing the resulting point agree, because lifts of a specified
base homeomorphism are unique once the value at one point is fixed.

A deck transformation preserving $L$ normalizes $\langle s\rangle$ and hence
centralizes $s$. Conversely a deck transformation centralizing $s$ preserves
$L$. Therefore the deck group of $L\to Z$ is exactly $C_\pi(s)$. Since $L\cong\R$,
this proves both the asserted isomorphism and injectivity of $\pi_1(Z)\to\pi$.
A finite-type real locus is a finite union of circles and open intervals after
its real punctures are removed. Each interval is homeomorphic to $\R$.
This argument is also the explicit reduction of Mochizuki's Theorem 3.6 and
Corollary 3.7 \cite{bMochizuki2003} to Grothendieck's formulation: his full
centralizer is $\pi_1(Z)\times C_2$, and intersecting it with $\pi$ gives the
centralizer used here.
\end{proof}

For $G=1$ the section set has one element and its centralizer is $\pi$; the
comparison is the bijection $\pi_0(U)\to\{*\}$. The map from individual points
is not injective. The hypothesis on the universal cover has content: on $S^2$
a nontrivial rotation has two fixed points while its lift extension, with kernel
$1$, has one section; the antipodal involution has no fixed point and still has
one section. These examples explain exactly why the same formula does not extend
to that excluded surface.

\begin{claim}{B.5}{Proved}{Recovering the incidence and cyclic order of real punctures}
For a hyperbolic $U=X\setminus S$ with anti-involution $\sigma$, its discrete
extension together with its indexed peripheral subgroups determines the following
finite graph. Its vertices are the section classes of B.4, its edges are the real
punctures, and each edge is attached to its two adjacent real half-branches.
Every punctured oval gives one cycle, allowing a loop or two parallel edges.
The graph recovers which punctures lie on each oval and their cyclic order up
to reversal. Unpunctured ovals are the vertices with infinite cyclic centralizer.
The claim uses the extension and peripheral incidence maps; a naked set of
section classes alone has no such incidence information.
\end{claim}
\begin{proof}
A puncture is real exactly when its $C_2$-stabilizer is all of $C_2$. For such a
puncture its peripheral extension is
\[
 E_i=\langle t,s\mid s^2=1,;sts^{-1}=t^{-1}\rangle,
 \qquad L_i=\langle t\rangle.
\]
Indeed the local anti-involution of the filled disc is reflection. On the
logarithmic covering $z=\exp(2\pi i w)$ of the punctured disc its two types of
order-two lifts are $w\mapsto-\overline w$ and
$w\mapsto1-\overline w$. Algebraically the reflections $t^k s$ have two
$L_i$-conjugacy classes, distinguished by the parity of $k$, since conjugation
by $t^m$ adds $2m$. The two inclusions of these classes into the section set of
$E$ are precisely the two real half-branches. They can have the same image.
Attach the edge $i$ using these two maps. Walking through the intervals of a
punctured circle gives a cycle whose edge labels are its punctures. Conversely
these cycles determine the partition of the punctures and their circular order,
with the orientation of each circle unchosen. B.4 identifies every remaining
vertex with an unpunctured oval. A nonreal pair of punctures has trivial
stabilizer and supplies no edge. Finally any two configurations on circles with
the same circular order are related by a homeomorphism: map successive closed
arcs monotonically, and extend over collars. Thus this recovery specifies the
topological position of the punctures, including the single-puncture and
single-interval case, without specifying their analytic positions.
\end{proof}

\subsection{Finite automorphisms and the withdrawn assertions}

\begin{claim}{B.6}{Proved}{The fundamental loop lemma of \S18}
Let $\pi=\pi_1(U)$ be hyperbolic and discrete, with its indexed peripheral
structure and orientation module. A finite-order orientation-preserving
automorphism of $\pi$ fixing a peripheral index and normalizing its subgroup
$L_i$ is the identity. Consequently the peripheral extension attached to a
nontrivial finite-order positive outer automorphism fixing $i$ cannot split.
This proves the lemma and corollary of scans 193--194 without adding realizability
to the hypotheses.
\end{claim}
\begin{proof}
The based Dehn--Nielsen--Baer correspondence identifies the given finite cyclic
group with a finite subgroup of the mapping class group with one additional
marked base point. Kerckhoff's realization theorem realizes this group by
homeomorphisms fixing that point and preserving the marked punctures. The
original paper explicitly treats cusps in \S IV, pp.~255--256, in addition to
the orientation-reversing convention on p.~235 \cite{bKerckhoff1983}; hence the
punctured case used here is included. The section determined by the fixed base
point identifies the induced genuine action on $\pi$ with the given one.

An orientation-preserving automorphism fixing $i$ acts as $+1$ on $L_i\cong\Z$,
so it fixes a generator $\ell$ of $L_i$. A nontrivial finite positive cyclic group
of plane lifts has one fixed point by B.1. Commutation with the nontrivial deck
transformation $\ell$ would give the second fixed point $\ell\widetilde x$,
a contradiction. Thus the automorphism is the identity. A splitting of the
peripheral extension gives a finite-order lift normalizing $L_i$; the result
just proved makes its outer class trivial.

The margin on scan 194 writes the \emph{filled} type as $(g,\nu)$ and the
acting group as type $(g,\nu+1)$. For $(g,\nu)=(0,2)$ or $(1,0)$ the acting
group is still a hyperbolic free group of rank two, so the same proof applies.
For a literal annulus, a positive automorphism fixing one of its two indexed
ends acts trivially on $\pi=\Z$. Forgetting the indices permits the positive
end interchange, which induces $-1$ on $\Z$ while normalizing its common
peripheral subgroup. For a literal unpunctured torus there is no peripheral
index, so the displayed loop lemma has no instance. These are the precise
comparisons required by the marginal remark, also reflected in A.13.
\end{proof}

\begin{claim}{B.7}{Refuted}{The unrestricted fixed-subgroup question, with its replacement}
The assertion $\pi^G=1$ for every finite $G\le\Aut_{\rm lac}(\pi)$ is false.
For a hyperbolic discrete $\pi$, a nontrivial finite genuine loop action has
fixed subgroup $1$ when its positive part is nontrivial. Its remaining possibility
is a reflection action of $C_2$; its fixed subgroup is $1$ or $\Z$, as in B.4.
For any finite-type faithful geometric action there are finitely many classes
of full and partial sections. This answers questions (c) and (e), scans 197--198.
\end{claim}
\begin{proof}
A finite subgroup of $\Aut_{\rm lac}(\pi)$ intersects $\Inn(\pi)$ trivially:
$\pi$ is torsion-free and centre-free. Its outer image is therefore finite and
is realized by Kerckhoff. The canonical centre-free extension identifies the
given genuine action with a section of its geometric lift extension. B.3 and
B.4 now give the complete list of fixed subgroups. An explicit counterexample
is the double of an orientable one-holed torus, a genus-two surface, with the
reflection interchanging the two halves. Its fixed oval is essential and a
base point on it gives a genuine order-two automorphism with fixed subgroup
$\Z$.

For finiteness, compactify $U$ by its finitely many ends. Every nonidentity
positive periodic homeomorphism has isolated fixed points by B.1. The resulting
closed discrete subset of the compact surface is finite. A reflection has a
compact one-dimensional fixed submanifold, hence finitely many components;
removing finitely many punctures preserves finiteness. Apply B.3 or B.4 to each
of the finitely many subgroups of $G$. This proves the asserted finiteness.
For the abelian groups the genuine actions can be read directly: a nonidentity
positive finite-order element of $SL_2(\Z)$ fixes no nonzero lattice vector,
whereas a determinant $-1$ involution fixes a primitive rank-one sublattice.
For $\Z$, the nonidentity automorphism is $-1$ and its fixed subgroup is zero.
The trivial group of operators fixes all of $\pi$ in every type.
\end{proof}

\begin{claim}{B.8}{Refuted}{Unique maximal sections fail for reflections, already in type $(0,3)$}
The orientation-free extension of question (d), scan 198, and of the expectation
on scan 155 is false. Put
\[
 E=\langle s_0,s_1,s_2\mid s_j^2=1, (s_0s_1)^2=1, (s_1s_2)^3=1\rangle.
\]
There is a geometric extension $1\to\pi_{0,3}\to E\to S_3\times C_2\to1$.
The reflection subgroup $\langle s_1\rangle$ is contained in two distinct maximal
finite subgroups $\langle s_0,s_1\rangle\cong C_2^2$ and
$\langle s_1,s_2\rangle\cong S_3$, which are not $\pi_{0,3}$-conjugate.
The valid uniqueness statement is B.3, requiring a nontrivial positive part.
\end{claim}
\begin{proof}
Take a hyperbolic triangle with angles $\pi/2,\pi/3,0$ and reflect in its three
sides. Reflected triangles tile the plane; the local tilings at the two finite
vertices give exactly the displayed relations. Equivalently its group is the
amalgam $C_2^2*_{C_2}S_3$, with common subgroup generated by $s_1$.
Write $S_3=\langle r,a\mid r^3=a^2=1, ara=r^{-1}\rangle$, and let $\tau$ generate
the central $C_2$. Define the surjection
\[
 s_0\mapsto\tau,\qquad s_1\mapsto a\tau,\qquad s_2\mapsto ra\tau.
\]
It is injective on both vertex groups. A finite subgroup of this amalgam fixes
a vertex of its Bass--Serre tree: the finite orbit spans a finite tree, whose
centre is fixed (there are no edge inversions). Thus the kernel is torsion-free.
The character sending each $s_j$ to $-1$ factors through the exponent of $\tau$,
so the kernel preserves orientation. The group Euler characteristic is
$1/4+1/6-1/2=-1/12$, and the index is $12$; hence the kernel has Euler
characteristic $-1$. The stabilizer of the cusp is $\langle s_0,s_2\rangle$;
its image has order four. The kernel therefore has $12/4=3$ cusps, and its
quotient surface has genus zero. This constructs the asserted geometric extension.

Each vertex group is maximal finite, again by the tree action: a finite overgroup
would fix a vertex, and a vertex group of order greater than the edge group fixes
only its vertex. Their intersection is $\langle s_1\rangle$ by the normal form
for the amalgam. Their orders differ, and their generated group is all of the
infinite amalgam. This proves every claimed failure. The finite quotient, its
vertex and cusp images, and the Euler-characteristic calculation are checked in
\texttt{code/check\_part\_b.py}. Grothendieck's proof on scans 201--203 applies
when the finite overgroups are positive cyclic groups: they centralize their
intersection, and B.3 makes the kernel of that centralizer map trivial. A
reflection has a line of fixed points, so that argument supplies neither of
these requirements in the generalization just refuted.
\end{proof}

\begin{claim}{B.9}{Proved}{The valid normalizer assertion}
For a hyperbolic discrete group with loops, a nontrivial finite subgroup
$K\le\Aut^+_{\rm lac}(\pi)$ is cyclic and
\[
 N_{\Aut^+_{\rm lac}(\pi)}(K)=C_{\Aut^+_{\rm lac}(\pi)}(K).
\]
An orientation-reversing element of its full normalizer acts on $K$ by inversion.
Every finite positive overgroup of $K$ is cyclic. These are the valid assertions
within the passage labelled ``Ultra fort!'' on scans 198--199.
\end{claim}
\begin{proof}
Realize the outer image of $K$ and use its given section. B.3 makes $K$ cyclic.
The boundary rotation character of B.12 is injective on $K$ and is natural for
loop automorphisms with their action on the orientation module. A positive
normalizer therefore preserves every value of that character, so fixes every
element of $K$. A negative normalizer changes its sign, so inverts every element.
The same realization argument proves cyclicity of each finite positive overgroup.
This proof compares a normalizing automorphism on the canonical oriented boundary;
it makes no assertion that an infinite normalizer acts as a finite isometry group.
\end{proof}

\begin{claim}{B.10}{Refuted}{The withdrawn uniqueness of topological roots}
The lemma marked ``faux!'' on scan 200 is false even for the punctured torus:
two different finite-order homeomorphisms fixing the base point can have the
same nontrivial square and induce the same automorphism of $\pi_1$.
\end{claim}
\begin{proof}
On $U=\R^2/\Z^2\setminus\{0\}$ choose $p=(1/2,1/2)$ and let
$r(x,y)=(-y,x)$, $h=r^2$. Choose a small disc $D$ whose four $r$-translates are
disjoint and avoid both $0$ and $p$. Choose a nonidentity homeomorphism $f_0$
supported in $D$ and isotopic to the identity there relative to its boundary;
for example use a radial bump rotation in disc coordinates. Put
$f=f_0(hf_0h^{-1})$. Its two disjoint supports show that $fh=hf$; the same
supports show that $fr\ne rf$, since $r$ moves their union to the disjoint pair
$rD,r^3D$. Hence $v=frf^{-1}\ne r$, but $v^2=h$ and $v(p)=p$.
The isotopy of $f$ fixes $0$ and $p$, so $f_*=1$ on $\pi_1(U,p)$ and therefore
$v_*=r_*$. This verifies the exact data of the withdrawn statement. What survives
is uniqueness up to conjugacy by the identity component in the realization
comparison, not equality of the homeomorphisms.
\end{proof}

\begin{claim}{B.11}{Refuted}{Torsion in the centralizer need not be a subgroup}
The assertion following the normalizer assertion on scan 199, that the torsion
elements of the centralizer form one largest finite cyclic overgroup, is false
for the hyperbolic type $(1,1)$.
\end{claim}
\begin{proof}
Keep $U,p,r,h$ as in B.10. The essential simple curve
$c=\{x+y=1/2\pmod1\}$ avoids $0,p$ and is invariant under $h$.
The based mapping class $T_c$ commutes with $h$, since conjugation of a positive
Dehn twist by an orientation-preserving map is the positive twist about the
image curve. Thus the two genuine automorphisms induced by $r$ and
$T_crT_c^{-1}$ have order four and common square $h_*$; both centralize $h_*$.
On $H_1(U;\Z)$, with a consistent choice of twist sign, their matrices are
\[
 S=\begin{pmatrix}0&-1\\1&0\end{pmatrix},\qquad
 TST^{-1}=\begin{pmatrix}-2&-5\\1&2\end{pmatrix},\qquad
 T=\begin{pmatrix}2&1\\-1&0\end{pmatrix}.
\]
Both square to $-I$, while their product has trace $-6$ and determinant one.
Its eigenvalues are $-3\pm\sqrt8$, so it has infinite order. The corresponding
product of automorphisms has infinite order as well. These integral equalities
are checked by the supplied program. Consequently the centralizer's torsion
set is not closed under multiplication, and there is no finite subgroup
containing all its finite overgroups of $\langle h_*\rangle$. This coexists with
the exact equality of normalizer and centralizer in B.9.
\end{proof}

\subsection{The multiplier and the abelian surfaces}

\begin{claim}{B.12}{Proved}{An intrinsic discrete formula for the multiplier}
For a geometric finite extension of a hyperbolic discrete surface group,
$E$ with its peripheral structure and orientation module $T$ determines
$\mu_E$ on all torsion elements. On $E^+$ it is the rotation number of its
canonical oriented boundary action. More explicitly, let
$\widetilde{\operatorname{Homeo}}^+(S^1)\to\operatorname{Homeo}^+(S^1)$ be the
central extension by the positively oriented deck translation $z$. For
$u\in E^+$ of order $n$ and any lift $\widetilde u$, write
$\widetilde u^{,n}=z^k$. Then
\[
 \mu_E(u)=\frac{k}{n}\pmod{\Z}\quad\hbox{in }T\otimes\Q/\Z.
\]
For a nonidentity orientation-reversing torsion element we put
$\mu_E(u)$ equal to the unique nonzero element of order two. With this
convention the formula is natural under isomorphisms preserving the loops and
transporting $T$, satisfies $\mu(u^m)=m\mu(u)$, and has the same order as $u$.
This answers the discrete intrinsicity question of \S15, folios 77--79,
scans 155--160.
\end{claim}
\begin{proof}
Choose an oriented finite-area hyperbolic uniformization of $U$. Its boundary
at infinity is a circle. We explain the ambiguity in this choice. A marking
isomorphism of surface groups preserving their peripheral subgroups and the
orientation is induced by an orientation-preserving homeomorphism, by the
Dehn--Nielsen--Baer correspondence. We can choose that marking quasiconformal:
on a compact core smooth it, and on each cusp use standard cylindrical
coordinates and extend with bounded dilatation. Its lift extends to an oriented
boundary homeomorphism. This extension is uniquely determined by its
$\pi$-equivariance. Indeed it must carry the attracting and repelling endpoints
of every hyperbolic element to those of the corresponding element. These
endpoints are dense in the boundary circle of a finite-area surface group.
For completeness, density follows from its full limit set: a boundary arc
meets translates of a fundamental ideal polygon arbitrarily near infinity;
conjugating either of two hyperbolic elements with disjoint endpoints by the
corresponding translates gives endpoints in that arc. Thus two equivariant
boundary maps agree on a dense set and hence everywhere.

For each automorphism of $\pi$ preserving the peripheral structure,
Dehn--Nielsen--Baer provides the boundary map conjugating its deck action.
The preceding uniqueness makes composition exact and makes this circle,
together with its $E$-action, independent of the uniformization up to a unique
$\pi$-equivariant boundary isomorphism. The boundary orientation is recorded by
$T$, so changing orientation acts by $-1$. Realize the finite action and choose
an invariant conformal structure, using B.1's cyclic and dihedral local charts
on the quotient. Its hyperbolic metric makes a finite-order positive lift an
elliptic isometry. In the disc model centred at its fixed point it is the same
rotation on the tangent circle and on the circle at infinity. The displayed
formula therefore gives its local multiplier.

Replacing $\widetilde u$ by $z^j\widetilde u$ replaces $k$ by $k+nj$, so the
value is independent of the lift. An orientation-preserving finite-order circle
homeomorphism of order $n$ is conjugate to a rotation by $k/n$ with
$\gcd(k,n)=1$; this also follows by ordering the finite orbit on the circle
and extending a homeomorphism of one fundamental interval around that orbit.
This proves the order and power statements. A negative finite plane isometry
is a reflection, of order two, so its assigned value has exactly the required
power property and is unchanged by sign reversal.

The normalizer formula suggested in the margin of scan 159 is recovered as
follows. Drill the orbit of the fixed point of $u$. Above the new puncture the
positive peripheral normalizer is infinite cyclic, generated by a positively
oriented element $v$, with peripheral subgroup generated by $v^n$. A lift of
$u$ to this normalizer is $v^a$ modulo $\langle v^n\rangle$, and
$\mu(u)=a/n$. In a local coordinate this is the equality between the rotation
$z\mapsto e^{2\pi ia/n}z$ and monodromy in $z\mapsto z^n$. Harvey's Theorem 7
\cite{bHarvey1971} computes exactly this inverse-residue relation from the
elliptic generators of the orbifold epimorphism. His convention uses an inverse
deck transformation, accounting for the sign visible in his proof on
pp.~393--394. Our local-coordinate calculation also applies to cusps and to the
positive cyclic part of a dihedral stabilizer; B.1 supplies the reflection part.
\end{proof}

\begin{claim}{B.13}{Refuted}{Forgetting loops can change the multiplier}
An abstract extension, with its orientation line regarded as an additional
rank-one module, need not determine the multiplier after the peripheral
structure is forgotten. There are isomorphic extensions with kernel a free
group of rank two and quotient $C_3$ whose geometric multipliers disagree,
even after allowing either isomorphism between their orientation lines.
\end{claim}
\begin{proof}
Let $E_0=\langle u,v\mid u^3=v^3=1\rangle=C_3*C_3$. Realize it as the
positive orbifold group of a sphere with one cusp and two cone points of order
three; choose $u,v$ as its positive local generators. The epimorphisms
\[
 q_+(u)=q_+(v)=1,\qquad q_-(u)=1,\quad q_-(v)=-1
 \quad (\text{values in }\Z/3)
\]
are injective on every finite subgroup, because finite subgroups of a free
product are conjugate into its factors (the tree argument in B.8 applies).
Their kernels are torsion-free of index three and Euler characteristic $-1$.
The cusp generator is $(uv)^{-1}$. Its image under $q_+$ is nonzero, so there
is one cusp and the kernel is a once-punctured torus group. Its image under
$q_-$ is zero, so there are three cusps and the kernel is a pair-of-pants group.
Both have free rank two.

The automorphism $\psi(u)=u,\ \psi(v)=v^{-1}$ satisfies $q_-\psi=q_+$ and
is an isomorphism of the two abstract extensions. The positive local generator
values in each realization are $\mu(u)=\mu(v)=1/3$. Thus $\psi$ preserves the
value at $u$ but changes the value at $v$ to $-1/3$. Reversing the isomorphism
of orientation lines would instead break the equality at $u$. No orientation-line
identification repairs both. The isomorphism sends the cusp word $uv$ to
$uv^{-1}$, displaying precisely which peripheral comparison has been discarded.
The finite monodromy calculation is included in the computational record.
\end{proof}

\begin{claim}{B.14}{Proved}{Finite annulus and torus actions, and their multipliers}
Every finite topological action on the torus is topologically conjugate to an
affine action on $V/\Lambda$, where $\Lambda\cong\Z^2$ and
$V=\Lambda\otimes\R$. For a fixed finite operator group $G$ such an action is
specified by a homomorphism $A:G\to GL(\Lambda)$ and a cocycle
$b\in Z^1(G,V/\Lambda)$, subject to faithfulness of
$g\mapsto(x\mapsto A_gx+b_g)$. Translation conjugacy corresponds exactly to a
coboundary, and
\[
 H^1(G,V/\Lambda)\xrightarrow{\sim}H^2(G,\Lambda)
\]
is the connecting homomorphism. Linear changes of the lattice give the further
conjugacies. Every finite annulus action has a simultaneous normal form
\[
 (x,t)\longmapsto(\epsilon_g x+a_g,\eta_g t),
 \qquad x\in\R/\Z,\quad t\in\R,\quad
 \epsilon_g,\eta_g\in\{1,-1\}.
\]
These descriptions determine the multiplier of every finite-order element in
the respective lift extensions.
\end{claim}
\begin{proof}
We first justify simultaneous conformalization. B.1 gives genuine cyclic or
dihedral charts for every finite stabilizer. The quotient is a compact
surface orbifold after the ends are added; its interior special points have
cyclic charts, and its mirror edges and corners have reflection and dihedral
charts. Choose a smooth metric on the quotient away from these charts, choose
round cone and mirror metrics in them, and patch in collars using an invariant
partition of unity on their finite local covers. The local charts make the
patched lifts smooth. Pullback is a smooth invariant conformal structure on
the original surface. This construction proves the required passage from the
original topological action; it does not average a merely topological metric
and call the result Riemannian.

Uniformization of the torus gives $\C/\Lambda$. Every holomorphic automorphism
lifts to an entire map $f$ with $f(z+\lambda)-f(z)\in\Lambda$. Its derivative
is periodic and hence bounded on $\C$, so it is constant; $f(z)=az+c$.
An antiholomorphic automorphism is obtained by composing such a map with
complex conjugation on its source. Hence every element is real affine on the
same lattice. The action law is precisely
$b_{gh}=b_g+A_gb_h$. Averaging cochains over the finite group gives
$H^q(G,V)=0$ for $q>0$. The exact sequence
$0\to\Lambda\to V\to V/\Lambda\to0$ therefore gives the asserted isomorphism.
In concrete terms, lifts $v_g\in V$ of $b_g$ give the extension cocycle
$v_g+A_gv_h-v_{gh}\in\Lambda$. Conjugating by translation adds a coboundary.
Moreover any topological conjugacy between two affine actions can be replaced
by an affine conjugacy. A lift of the conjugacy has form $Bx+f(x)$, with
$B\in GL_2(\Z)$ and $f$ periodic. Average its equivariance equation over the
torus with Haar measure. Affine lattice automorphisms preserve that measure,
so $Bx+\int f$ satisfies the same equation modulo $\Lambda$.

For the annulus, extend the action to its two-point sphere compactification
and apply Kolev. In the resulting conformal sphere choose a M\"obius coordinate
sending the two ends to zero and infinity. The maps are then of the forms
$z\mapsto az$, $a/z$, $a\overline z$, or $a/\overline z$. On the logarithmic
radial coordinate they act affinely with slopes $\eta=\pm1$. A finite group
of affine transformations of the real line fixes the average of any orbit;
translation to this common fixed point removes every radial constant. The
angular coordinate gives the claimed $\epsilon x+a$. This proves the displayed
simultaneous normal form. The action on $\pi_1$ is $\epsilon_g$ and its
orientation sign is $\epsilon_g\eta_g$.

For the promised computation on the torus, an element of the lift extension
is $x\mapsto Ax+v$. It has finite order exactly when $A$ has finite order,
say $n$, and $(1+A+\cdots+A^{n-1})v=0$. Averaging its orbit gives a fixed point,
and translation to that point conjugates it to $A$. A nonidentity positive
finite-order matrix has trace $-2,-1,0,1$, giving orders $2,3,4,6$, respectively.
For order two its multiplier is $1/2$. For the other orders it is $+1/n$ or
$-1/n$ according to the sign of $\det(w,Aw)$ for any nonzero $w$; this sign
is independent of $w$, as is seen after averaging an inner product and writing
$A$ as a rotation. A negative finite-order matrix is a reflection and receives
$1/2$. A lift with $A=I$ has finite order only when it is the identity.
For the annulus a lift is $(x,t)\mapsto(\epsilon x+a+k,\eta t)$. With
$\epsilon=-1$ its square is the identity; with $\epsilon=1$ finite order
requires $a+k=0$, leaving the identity or $t\mapsto-t$. Thus every nonidentity
torsion lift on an annulus has order two and multiplier $1/2$.
\end{proof}

\begin{claim}{B.15}{Refuted}{The ``only counterexample'' of scan 179}
The assertion that a faithful homotopically trivial finite annulus action
failing to be a translation action must be the single anti-involution
$z\mapsto1/\overline z$ is false. For the two non-simply-connected abelian
types, the exact replacement is: homotopically trivial torus actions are
translation actions; homotopically trivial annulus actions are the finite
subgroups of $SO(2)\times C_2$ in the normal form of B.14. Those contained in
$SO(2)$ are translations, while those with nontrivial second projection include
two different types of anti-involution and groups of arbitrarily large order.
\end{claim}
\begin{proof}
On the torus, homotopic triviality is $A_g=I$ for every $g$, so B.14 proves the
translation statement and the implication doubted on scans 180--181.
On the annulus it is $\epsilon_g=1$, which gives the asserted subgroup.
For each $n\ge1$ the maps
\[
 r(z)=e^{2\pi i/n}z,\qquad s(z)=1/\overline z
\]
generate a faithful $C_n\times C_2$ action, and every element induces the
identity on $\pi_1(\C^*)$. The fixed circle of $s$ excludes a translation
realization: a nonidentity translation in any group is fixed-point free.
In addition $s_-(z)=-1/\overline z$ is an involution inducing $+1$ on $\pi_1$,
but it has no fixed point, since its fixed-point equation is $|z|^2=-1$.
It cannot be conjugate to $s$, and it cannot be a translation of a connected
topological group: every such translation is isotopic to the identity along
a path in that group and hence preserves orientation. These formulas give the required full
counterexamples, not merely a gap in the argument.

More explicitly let the rotation subgroup have order $n$ and generator
$(x,t)\mapsto(x+1/n,t)$. A second-coset generator can be chosen as
$(x,t)\mapsto(x+m/(2n),-t)$, with $m=0$ or $1$ after multiplying by a rotation.
Its square is the $m$th rotation. This exhausts the finite subgroups projecting
onto $C_2$, and includes both the reflection and glide cases. For the omitted
simply connected surfaces, finite plane actions are the orthogonal actions of
B.1 and are all homotopically trivial; positive sphere actions are likewise
homotopic to the identity. Nontrivial rotations fixing a point already refute
any extension of the translation formulation to those types.
\end{proof}

\subsection{Anti-involutions and the later tube construction}

\begin{claim}{B.16}{Proved}{Complete topological classification of anti-involutions}
Conjugacy classes of anti-involutions of a closed oriented genus-$g$ surface
are classified by the orientability of the quotient and its number $b$ of
boundary components. The possibilities are
\[
 \begin{array}{ll}
 \text{nonorientable quotient:}&0\le b\le g,\quad
                  h=g+1-b\ge1\text{ crosscaps};\\
 \text{orientable quotient:}&b\ge1,\quad g=2h+b-1,\quad h\ge0.
 \end{array}
\]
There are $(g+1)+\lfloor(g+2)/2\rfloor=\lfloor(3g+4)/2\rfloor$ classes.
For $\nu$ unlabelled invariant punctures the additional data are an integer
$r\ge0$ and an unordered list $(q_1,\ldots,q_b)$ of nonnegative integers with
$2r+\sum q_i=\nu$: $r$ is the number of interior punctures of the quotient,
and $q_i$ counts punctures on its $i$th boundary circle.
\end{claim}
\begin{proof}
By B.1 the quotient $Y=X/\sigma$ is a compact surface with boundary
$X^\sigma$. The complement of the real locus covers its interior with degree
two; the Euler characteristic of the boundary is zero. Hence
$\chi(X)=2\chi(Y)$. If $Y$ is nonorientable with $h$ crosscaps,
$\chi(Y)=2-h-b$, yielding $g=h+b-1$; if orientable of genus $h$ it yields
$g=2h+b-1$. In the latter case $b$ must be positive, since the oriented double
of a closed orientable quotient would be disconnected. Conversely form the
orientation double over the interior of each allowed $Y$ and join the two
local sides at its boundary. The result is a connected closed orientable surface,
with the deck interchange extended to reflection along the boundary; its genus
is given by the displayed formula. The quotient recovers $Y$, so the classification
of compact surfaces proves both existence and uniqueness of these conjugacy classes.
This also proves the complete list underlying Gross--Harris, Proposition 3.1
\cite{bGrossHarris1981}, including the two corrections in the margins of scan 439.
Grothendieck's nonorientable parameter $\gamma$ is $h-1$; thus his formula
$g=\gamma+b$ agrees exactly with ours.

An anti-involution of a punctured finite-type surface extends to its end
compactification. Punctures in its quotient lie either in the interior,
with two-point inverse image, or on the boundary, with one-point inverse image.
Any two finite interior configurations of the same cardinality are related by
an isotopy supported away from the boundary. Boundary components may be
permuted, and any two unlabelled configurations of the same cardinality on one
circle are related by a circle homeomorphism, extended over a collar. These
facts prove completeness of the punctured invariants. With individually
labelled punctures one must instead retain their assignments and their circular
orders; B.5 provides the corresponding incidence reconstruction.
\end{proof}

\begin{claim}{B.17}{Proved}{Real moduli components, with the group and marking convention fixed}
Let $\mathcal M_{g,\nu}$ mean the real moduli stack of smooth proper
geometrically connected genus-$g$ curves with an \emph{unlabelled}
degree-$\nu$ reduced puncture divisor. We retain the stack and its
real forms also in the four unstable types. Its complex
mapping class group is $\mathcal T^+_{g,\nu}$, allowing permutations of the
punctures. Write $p_b(N)$ for the number of partitions of $N$ into at most
$b$ parts, with $p_0(0)=1$ and $p_0(N)=0$ for $N>0$. Then
\[
 \#\pi_0\bigl(\mathcal M_{g,\nu}(\R)\bigr)
 =\sum_{b=0}^{g}\ \sum_{r=0}^{\lfloor\nu/2\rfloor}p_b(\nu-2r)
  +\sum_{\substack{1\le b\le g+1\\ b\equiv g+1\ (2)}}
       \ \sum_{r=0}^{\lfloor\nu/2\rfloor}p_b(\nu-2r).
\]
For closed hyperbolic surfaces this is $\lfloor(3g+4)/2\rfloor$.
In the four unstable types $(0,0)$, $(0,1)$, $(0,2)$ and $(1,0)$ the counts are,
respectively, $2,1,3,3$. In the last case no origin is prescribed.
Distinct topological anti-involution classes remain distinct in the extended
mapping class group in these hyperbolic types. The last assertion fails for
$(0,0)$, $(0,2)$, and $(1,0)$.
\end{claim}
\begin{proof}
C.5 and C.6, finite-action realization and uniqueness up to the identity component, identify
conjugacy classes of anti-involutions on the fixed punctured surface with
conjugacy classes of orientation-reversing involutions in
$\mathcal T_{g,\nu}$. One can use positive conjugacy here: an orientation-reversing
conjugator can be multiplied by the target involution to become positive.
Mochizuki, Corollary 3.8, identifies these conjugacy classes with the connected
components of the real moduli \emph{stack} \cite{bMochizuki2003}. His Example 3.3
uses the full modular group of a punctured surface, which is the unlabelled
convention just specified. B.16 supplies the list. For fixed $b,r$, unordered
nonnegative lists of length $b$ summing to $\nu-2r$ are exactly the partitions
counted by $p_b(\nu-2r)$. Summing over the two quotient families proves the
formula. The supplied code enumerates these numbers for $2g+\nu\le6$ and checks
the closed formula for $0\le g\le20$; the proof holds for all $g$.

We verify the unstable component counts without using mapping-class
injectivity. A real genus-zero curve is a split or nonsplit real conic;
analytically these are the sphere with $z\mapsto\overline z$ or with
$z\mapsto-1/\overline z$. To see these are the two possibilities,
conformalize its anti-involution as in B.14 and apply the two quotient
types in B.16; uniformization identifies the resulting conformal
double with the sphere. A real point exists exactly in the first type.
There is no complex moduli parameter in genus zero. With one puncture
the point is real, and $PGL_2(\R)$ acts transitively on the real
projective line, giving one component. With two unlabelled punctures
the possibilities are two real points on the split conic, one
conjugate pair on the split conic, or one conjugate pair on the
nonsplit conic. The first is transitive under $PGL_2(\R)$, the second
under its action on the upper half-plane, and the third under the
rotations commuting with the antipodal map. These give three
components. With no puncture there are the two conic components.

For an unpointed genus-one curve, B.14 and B.16 give the three
anti-involution types with zero, one, or two real ovals. Each type
has a connected parameter space of complex structures. Explicit
representatives are
\[
 \begin{array}{c|c|c}
 b&\Lambda_t&\sigma(z)\\\hline
 0&\Z+it\Z&\overline z+\tfrac12\\
 1&\Z+(\tfrac12+it)\Z&\overline z\\
 2&\Z+it\Z&\overline z
 \end{array}\qquad(t>0).
\]
To obtain this list from an invariant lattice, rotate a primitive
real period to $1$. Conjugation of a second period forces its real
part modulo $\Z$ to be $0$ or $1/2$. The translation cocycle of
B.14 is then removed by a change of origin, except for the
nonzero half-period class in the rectangular case. That class gives
the first row. Thus every real genus-one curve occurs in one of
these connected families. The number of real ovals is locally
constant in a smooth proper real family: local trivialization can
be chosen invariant under the involution by an invariant horizontal
metric and its path lifting. Distinct rows consequently belong to
different components. This proves the number three, with no origin
added to the moduli problem. Direct substitution of the four
unstable types into the displayed partition formula gives exactly
these counts.

The mapping-class exceptions can be checked directly. On $S^2$ reflection and the antipodal
map have different fixed sets, but both represent the orientation-reversing
mapping class. On $\R^2/\Z^2$ the maps $(x,y)\mapsto(x,-y)$ and
$(x,y)\mapsto(x+1/2,-y)$ have respectively two fixed circles and no fixed
points, but both induce $\operatorname{diag}(1,-1)$ and hence the same mapping
class. On the annulus, the reflection and glide anti-involutions of B.15 also have
the same mapping class: the path $(x,t)\mapsto(x+s/2,-t)$ joins them through
homeomorphisms preserving the two ends as a set. These examples establish the exact failure of the unrestricted
injectivity hoped for on scan 440. For labelled moduli spaces or a coarse
moduli space, the quotient and its real locus have to be specified separately;
the displayed formula is for the stack and puncture convention stated above.
\end{proof}

\begin{claim}{B.18}{Proved}{The two centralizer corollaries in both quotient families}
Let $g\ge2$, and let an anti-involution of $X_g$ have $b>0$ ovals, with outer
class $\sigma$. There are exactly $b$ classes of order-two genuine
loop automorphisms lifting $\sigma$, under conjugacy by $\pi_g$.
For each lift $u$, $\pi_g^u\cong\Z$. The positive centralizer of $\sigma$ acts
transitively on those $b$ classes. More precisely, writing
$A=\Aut_{\rm lac}(\pi_g)$ and $O=\Out_{\rm lac}(\pi_g)$, there is an exact sequence
\[
 1\longrightarrow\pi_g^u\longrightarrow C_{A^+}(u)
 \longrightarrow\operatorname{Stab}_{C_{O^+}(\sigma)}([u])\longrightarrow1,
\]
and the last group has index $b$ in $C_{O^+}(\sigma)$. The corresponding full
centralizer image also has index $b$. This proves Corollaries 1 and 2,
folios 225--226, scans 445--448, for both orientable and nonorientable quotients.
\end{claim}
\begin{proof}
B.4 proves the count and the fixed subgroup. Any permutation of the boundary
components of a compact connected quotient surface is realized by a homeomorphism:
join two components by an embedded arc, interchange them in a neighbourhood of
that arc and the two collars, and use transpositions. Each quotient homeomorphism
lifts to the oriented double; among its two lifts choose the orientation-preserving
one. It commutes with the reflection. Consequently the positive centralizer is
transitive on the ovals, and hence on the section classes.

For exactness, the centre of $\pi_g$ is trivial, so $\Inn(\pi_g)$ is identified
with $\pi_g$. Its intersection with $C_A(u)$ is exactly $\pi_g^u$.
An outer class $\alpha$ centralizing $\sigma$ preserves $[u]$ precisely when
one, hence every, lift $\widetilde\alpha$ satisfies
$\widetilde\alpha u\widetilde\alpha^{-1}=vuv^{-1}$ for an inner automorphism
$v$. Then $v^{-1}\widetilde\alpha$ centralizes $u$, proving surjectivity onto
the displayed stabilizer. Orbit--stabilizer gives index $b$ in both versions.
For $b=0$ there are no order-two genuine lifts and these assertions about a
chosen lift have no instance. With punctures the same exact sequence holds,
but the orbit size records the boundary decorations from B.16; it need not be
the total number of ovals when different ovals carry different puncture data.
\end{proof}

\begin{claim}{B.19}{Proved}{The tube twists, their maximality property, and their normalizer}
For $g\ge2$, the elements $\rho_j$ of \S33, scans 453--470, are full Dehn twists
about a cut system $c_1,\ldots,c_g$: disjoint nonseparating curves with connected
complement. They generate $H\cong\Z^g$. This is not a maximal free abelian
subgroup of $\mathcal T_g$: it is contained in a multitwist subgroup of rank
$3g-3$. It has the precise maximality property of being a full multitwist group
of largest rank whose supporting multicurve has connected complement.
Furthermore
\[
 N_{\mathcal T_g}(H)=\operatorname{Stab}_{\mathcal T_g}(\{c_1,\ldots,c_g\}),
 \qquad N_{\mathcal T_g}(H)/C_{\mathcal T_g}(H)\cong S_g\times C_2.
\]
The action on $H$ permutes its basis and multiplies every basis vector by the
same orientation sign. Cutting and capping gives an exact sequence
\[
 1\longrightarrow H\longrightarrow N_{\mathcal T_g}(H)
 \longrightarrow\mathcal T^{\rm pairs}_{0,2g}\longrightarrow1,
\]
where the last extended sphere mapping class group preserves the partition of
its $2g$ punctures into the $g$ unordered gluing pairs.
\end{claim}
\begin{proof}
In the coordinates of scan 461 the map is
$\rho_j(z,t)=(-z e^{\pi i t},t)$ on $S^1\times[-1,1]$ and is the identity outside
the tube. Its angular displacement increases from $0$ to $2\pi$, so its relative
boundary mapping class is one full Dehn twist, with sign fixed by that coordinate
orientation. The complement of the tubes in his construction is a sphere with
$2g$ boundary circles, establishing the cut-system assertion.

Twists in disjoint collars commute. They are independent: cut the surface along
the $c_j$, choose transverse curves dual to these cuts, and measure the integer
winding of the resulting arcs when the collars are glued back. A product
$\prod T_{c_j}^{n_j}$ changes the $j$th such winding by $n_j$; an isotopy cannot
change these relative-end winding differences. This proves $H\cong\Z^g$.
Complete the $2g$-holed sphere to a pants decomposition by adding $2g-3$
essential disjoint curves. The same winding argument gives a rank-$3g-3$
multitwist group containing $H$. In contrast, cutting $k$ disjoint curves while
keeping the complement connected reduces the genus to $g-k$, so $k\le g$.
This proves the stated maximality among connected-complement multitwist systems.

For the normalizer, conjugation satisfies
$fT_cf^{-1}=T_{f(c)}^{\epsilon(f)}$. A single twist can equal a product of twists
about the $c_j$ only by being one of these twists to the corresponding power.
Here is the needed support argument. Such a twist commutes with every $T_{c_j}$.
The elementary intersection formula
$i(T_d(c),c)=i(c,d)^2$ for a single twist forces $d$ to be disjoint from each
$c_j$; otherwise its conjugation changes $T_{c_j}$. Add $d$ to the disjoint
system if it is not already there and apply the preceding independent-winding
argument. The resulting nontrivial relation is impossible. Thus a normalizer
permutes precisely the $c_j$, proving the first equality. Conversely every such
permutation of curves normalizes $H$. Every permutation of the handles is
realized positively, and a reflection of the standard cut-system model reverses
the global orientation while preserving each curve. This gives the full image
$S_g\times C_2$. Its kernel is exactly the centralizer.

A representative preserving the multicurve can be isotoped to preserve its
collars. Its restriction to the cut surface, after capping each boundary to a
puncture, preserves the gluing pairs. Conversely a homeomorphism of the capped
sphere preserving those pairs can be made radial in small puncture discs and
then extended over the paired collars. Two such extensions differ by integer
collar twists. An isotopy on the capped sphere lifts after allowing exactly
these twists, by the isotopy extension theorem on collars. This proves both
surjectivity and the stated kernel, and types the comparison with the
complement suggested on scan 470.
\end{proof}

\begin{claim}{B.20}{Refuted}{Presentation (52) and its exact correction}
For $g\ge2$, presentation (52) on scan 469 does not present the stated
$\Z^g$-by-finite group. Put $r=\varepsilon_0\varepsilon_1$ and
$\rho_j=r^j\tau'\tau r^{-j}$, $0\le j<g$. The printed presentation gives
\[
 F_g\rtimes\bigl(D_{4g}\times\langle\tau\rangle\bigr),
\]
where $D_{4g}=\langle r,\varepsilon_1\mid r^{2g}=\varepsilon_1^2=1,
\varepsilon_1r\varepsilon_1=r^{-1}\rangle$ permutes the free generators through
its action on $\Z/g$, and $\tau$ inverts each free generator.
Adding the relations $[\rho_i,\rho_j]=1$ for $i<j$ gives exactly
$\Z^g\rtimes(D_{4g}\times C_2)$ with the manuscript's specified action.
\end{claim}
\begin{proof}
The relations say that $\tau$ commutes with $D_{4g}$, that $\tau'$ commutes
with $\varepsilon_1$ and $r^g$, and that both $\tau,\tau'$ are involutions.
Set $\rho_0=\tau'\tau$. Then $\tau\rho_0\tau=\rho_0^{-1}$,
$\varepsilon_1\rho_0\varepsilon_1=\rho_0$, and
$r^g\rho_0r^{-g}=\rho_0$. Its conjugates are indexed by the $g$ cosets of
$\langle\varepsilon_1,r^g\rangle$ in $D_{4g}$, with no further relation among
them. To check the last assertion directly, form the displayed semidirect
product with a free basis $\rho_0,\ldots,\rho_{g-1}$ and define
$\tau'=\rho_0\tau$ in it. All printed relations hold. The reverse map sends
these free generators to $r^j\tau'\tau r^{-j}$; the two maps agree inversely
on all generators. This is a proof by explicit mutually inverse homomorphisms.
Abelianizing only the free factor adds exactly the stated commutators and gives
the corrected group.

There is also a finite certificate. For $g=2$ set
$\varepsilon_0=(13)$, $\varepsilon_1=\tau=1$, $\tau'=(12)$ in $S_3$.
All nine displayed relators evaluate to $1$, but
$\rho_0=(12)$ and $\rho_1=(23)$ do not commute. The program checks every one
of these evaluations. Relations (33), scan 454, already assert commutation
for the geometric twists; they are exactly the missing input to (52).
\end{proof}

\begin{claim}{B.21}{Proved}{The one-holed-torus subgroups and the proved part of the analogy}
For $g\ge2$, choose a one-holed torus $W_j$ containing $c_j$ and a transverse curve $d_j$
with $i(c_j,d_j)=1$. Its group fixing its boundary pointwise embeds by extension
by the identity into $\mathcal T_g^+$. It is
\[
 B_3=\langle A_j,B_j\mid A_jB_jA_j=B_jA_jB_j\rangle,
\]
and capping the boundary gives
\[
 1\longrightarrow\langle T_{\partial W_j}\rangle\cong\Z
 \longrightarrow B_3\longrightarrow SL_2(\Z)\longrightarrow1,
 \qquad (A_jB_j)^6=T_{\partial W_j}^{\pm1}.
\]
This extension has no section. For disjoint chosen $W_j$, their supported
subgroups commute and, together with $H$, preserve their boundary multicurve;
they do not generate $\mathcal T_g^+$. On explicit generators the tube subgroup
and these two twists act by
\[
 \begin{array}{lll}
 \rho_j:&a_j\mapsto a_j,&b_j\mapsto b_ja_j,\\
 B_j:&a_j\mapsto a_jb_j^{-1},&b_j\mapsto b_j,
 \end{array}
\]
fixing the other handle generators of
$\pi_g=\langle a_j,b_j\mid\prod_j[a_j,b_j]=1\rangle$.
\end{claim}
\begin{proof}
Choose the based handle loops so that their commutator is the boundary loop.
The displayed substitutions preserve $[a_j,b_j]$ exactly, and are the effect
of crossing the appropriate twist collar. Twists in curves meeting once obey
the braid relation: in these generators substitution gives $ABA=BAB$.
The sixfold product $(AB)^6$ conjugates both generators by
$[a,b]^{-1}$, which is the boundary twist with this collar convention.
Both word equalities are checked by free reduction in
\texttt{code/free\_words.py}; the stated substitutions also give their direct
verification. Filling the boundary gives the standard two elementary matrices
generating $SL_2(\Z)$. Their action on primitive slopes, reduced by the Euclidean
algorithm, shows that the filled mapping class group is exactly $SL_2(\Z)$.
The kernel of filling is the boundary winding: cut along two based generating
arcs, isotope a kernel representative to fix them, and read its remaining
integer twist in the boundary collar. The kernel is therefore the infinite
cyclic group generated by the boundary twist. The braid group on three strands
has the same central extension: with $x=ABA,y=AB$ it is
$\langle x,y\mid x^2=y^3\rangle$, and its quotient by $y^6$ is $SL_2(\Z)$.
The comparison sends the kernel generator to the boundary twist, so it is an
isomorphism of these extensions.

Extension by the identity is injective because an isotopy to the identity on
the closed surface would fix the essential boundary and, after collar adjustment,
give an isotopy on the one-holed torus; a nonzero boundary twist remains
nontrivial by the winding argument in B.19. The braid group is torsion-free:
$\langle x,y\mid x^2=y^3\rangle$ is the amalgam of two infinite cyclic groups
along their subgroups of indices two and three; a finite subgroup fixes a vertex
in its tree and is conjugate into an infinite cyclic group. A section from
$SL_2(\Z)$ would inject its order-four element, contradicting this fact.

Disjoint supports commute. Each supported homeomorphism fixes the chosen
boundary curves, and so does every generator of $H$. An essential curve crossing
one of those boundaries has a twist which does not preserve it, by the
intersection formula in B.19. Hence the generated subgroup is proper.
The exact analogy proved here is the homomorphism to the usual rank-two
integral representation and its central boundary-winding extension, with braid
and disjoint-support relations. For overlapping choices of the strips one
must specify their intersections before computing the remaining relations;
no unspecified choice of strips supplies additional relations by itself.
\end{proof}

\subsection{Computations and the limits of the requested table}

\begin{claim}{B.22}{Refuted}{A finite subgroup alone does not index the requested abelian table}
For all types in Part B(f), a table indexed only by a finite subgroup of the
mapping class group cannot assign unique splitting numbers or $\Phi$.
The type $(1,0)$ and the single subgroup
$\langle\operatorname{diag}(1,-1)\rangle$ already have two different answers.
The additional extension datum is precisely the class in $H^2(G,\pi)$ in the
abelian cases of B.14.
\end{claim}
\begin{proof}
On the torus compare
$\sigma_0(x,y)=(x,-y)$ and
$\sigma_1(x,y)=(x+1/2,-y)$. Both are faithful involutions with the same
homomorphism to $GL_2(\Z)$, but the first fixes two circles and the second fixes
nothing. The corresponding extensions, with commuting lattice generators
$a,b$, have relations
\[
 \widetilde\sigma a\widetilde\sigma^{-1}=a,\qquad
 \widetilde\sigma b\widetilde\sigma^{-1}=b^{-1},\qquad
 \widetilde\sigma^2=a^\epsilon,\quad \epsilon=0\text{ or }1.
\]
The square of $a^mb^n\widetilde\sigma$ is $a^{2m+\epsilon}$. For $\epsilon=1$
there is no order-two lift. For $\epsilon=0$ such lifts require $m=0$,
and conjugation by $b$ changes $n$ by two. Thus there are exactly two section
classes in the first case and zero in the second. These are the two parity
classes of the extension cocycle. B.14 proves that retaining the extension
class gives the correct affine classification, with all translation conjugacies
accounted for. Thus the counterexample identifies the missing index and its
exact replacement.
\end{proof}

\begin{claim}{B.23}{Proved}{Complete finite-subgroup table in type $(0,3)$}
The ten conjugacy classes in the extended group
$\mathcal T_{0,3}\cong S_3\times C_2$, and their discrete section invariants,
are the following. In the table $F$ counts full-section classes, $P$ counts all
nontrivial partial-section classes with their subgroup images distinguished,
$M=|\Phi|$ counts maximal partial-section classes, and $O$ counts the group
orbits on $\Phi$.
\[
\begin{array}{lrrrrr}
 G&|G|&F&P&M&O\\\hline
 1&1&1&0&0&0\\
 C_2\text{ (positive)}&2&1&1&1&1\\
 C_2\text{ (horizontal mirror)}&2&3&3&3&3\\
 C_2\text{ (vertical mirror)}&2&1&1&1&1\\
 C_3&3&2&2&2&2\\
 C_2^2&4&1&6&3&2\\
 C_6&6&0&5&5&2\\
 S_3\text{ (positive)}&6&0&5&5&2\\
 S_3\text{ (mixed)}&6&2&7&2&2\\
 S_3\times C_2&12&0&16&5&2
\end{array}
\]
The attached CSV gives every maximal stabilizer; the JSON table gives every
partial section and all of its multiplier values. In particular it checks the
reflection failure of maximal-section uniqueness from B.8 inside this whole list.
\end{claim}
\begin{proof}
Place the three punctures at the vertices of an equatorial equilateral triangle.
The positive rotation group of that triangle is $S_3$: it is generated by the
$120$-degree vertical rotation $r$ and a half-turn $a$ around an axis through a
puncture. Reflection $\tau$ in the equator commutes with both. The mapping class
group map to permutations and orientation is an isomorphism. For its kernel,
fix three arcs forming a tree through the punctures by isotopy; its complement
is a disc, and Alexander's disc isotopy makes the remaining map isotopic to the
identity. This proves completeness of the ambient finite group and its stated
realization.

The program represents this group exactly by pairs $(p,e)\in S_3\times C_2$.
Starting with the trivial subgroup it repeatedly adjoins every group element
and takes closure, until no new subgroup occurs. Every subgroup is generated
by its elements, so this terminating procedure is exhaustive. It finds sixteen
subgroups and ten conjugacy classes, represented in the displayed order by
\[
 1,\ \langle a\rangle,\ \langle\tau\rangle,\
 \langle a\tau\rangle,\ \langle r\rangle,\ \langle a,\tau\rangle,\
 \langle r,\tau\rangle,\ \langle r,a\rangle,\
 \langle r,a\tau\rangle,\ \langle r,a,\tau\rangle.
\]
Here is the independent geometric input to its section calculation. Denote the
poles by $N,S$, and the three surviving opposite equatorial axis points by
$Q_i$; the other three axis points are the removed punctures. The fixed set of
$\langle r\rangle$ is $\{N,S\}$, and that of a positive half-turn is its single
surviving $Q_i$. The horizontal mirror has three open intervals, each containing
one $Q_i$. A vertical mirror has one real line after its single puncture is
removed, and this line contains $N,S,Q_i$. Intersect these fixed sets to find the
fixed sets of every subgroup. B.3 and B.4 turn their components into exactly
its section classes. A partial section on a mirror interval is contained in a
point stabilizer precisely when that point lies on the interval; this determines
all inclusions and hence all maximal classes. The code records and checks these
inclusions and the conjugation action of every element.

Every reflection and positive half-turn has multiplier $1/2$. For $r$, the
values at $N,S$ are respectively $1/3,-1/3$, and its powers give all remaining
values. This computes all entries in the attached JSON. The geometry also
verifies that all the centralizers occurring for nontrivial sections here are
trivial: every mirror component is a line, and every positive section fixes an
isolated point. The absence of a common finite overgroup for the two vertex
stabilizers along a vertical mirror is retained in the inclusion data.
\end{proof}

\begin{claim}{B.24}{Proved}{Completed finite-subgroup catalogue and exact indexing}
Every finite-subgroup conjugacy class in the requested range
$2g+\nu\le6$ has now been enumerated.  For the hyperbolic types,
the numbers of classes in the extended mapping class group are
\[
\begin{array}{c|r|l}
 (g,\nu)&\#\text{ classes}&\text{proofs and executed tables}\\\hline
 (0,3),(0,4),(0,5),(0,6)&10,30,20,57&B.23,B.25,B.26\\
 (1,1),(1,2),(1,3),(1,4)&13,40,38,115&B.25,B.27\\
 (2,0),(2,1),(2,2)&92,32,121&B.28,B.29\\
 (3,0)&292&B.33,B.37\text{--}B.41.
\end{array}                                             \tag{B.24.1}
\]
Each table includes the full and partial section classes, multipliers,
centralizers, extension incidences, maximal section classes $\Phi$,
and their group orbits and stabilizers.  In the four exceptional
abelian types, the finite mapping-class subgroup counts are
respectively $2,2,5,13$ for sphere, plane, annulus and torus.  Their
injectively realized finite-action counts are $3,2,7,17$; the exact
covering extension is retained as part of the section index in B.36,
as required by the explicit failures of extension-independent indexing
in B.22.  The families with a nontrivial kernel on passage to mapping
classes are described in B.14.
The exact tube words, stabilizer action, strip relations and arbitrary
transversal obstruction formerly required here are proved and executed
in B.30--B.32.
\end{claim}
\begin{proof}
The first three rows of (B.24.1) are precisely the complete catalogues
cited there.  We spell out the closed genus-three reconciliation to
record the coverage and the absence of double counting.
The positive portion has $55$ classes by B.33.  For a mixed finite
action let $H$ be its positive subgroup.  The positive quotient has
genus zero, one, two or three, by the positive classification.
The genus-zero case gives either a mirror disc ($145$ classes,
B.37) or a projective plane ($11$ classes, B.39).  The genus-one case
gives an annulus or M\"obius band ($57$ classes, B.40), or a Klein
bottle ($6$ classes, B.41).  The genus-two case is the free positive
$C_2$ action and gives $12$ mixed classes in B.38.  The genus-three
case has trivial positive part and gives its six negative involution
classes in B.38.  Thus the number of mixed classes is
\[
                   145+11+57+6+12+6=237,
                     \qquad55+237=292.                \tag{B.24.2}
\]
Every conjugacy descends to homeomorphisms of both the positive
quotient and the full quotient orbifold.  The positive quotient genus,
and the full quotient's boundary and orientability, separate the
listed portions.  Within each portion the cited proofs calculate
exactly the geometric kernel-normalizer orbits, so these counts are
irreducible counts of surface-action conjugacy classes.

B.34 supplies an independent hyperelliptic construction and is
counted through its explicit overlaps, not added a second time.
All $79$ of its rows are identified by monodromy isomorphisms:
$18$ with positive B.33 rows, $57$ with B.37 rows, and four with B.39
rows.  The executed program
\path{reconcile_genus3_catalogue.py} records these $79$ maps and
checks the corresponding section counts.  Its unified
\path{type_3_0_complete_catalogue.csv} has $292$ different
class identifiers, and \path{type_3_0_complete_sections.json}
contains the full finite group and section model for every row.
The exceptional replacement in B.36 is complete with its stated
realization index; it retains, in particular, the sphere antipodal
action's algebraic splitting without identifying it with a geometric
fixed point.  This proves the claimed coverage with exactly the
comparisons and indices used by the computations.
\end{proof}
